% !TEX root = ../matan.tex

\section{Степенные ряды и связанные факты}

\subsection{Степенные ряды в комплексной плоскости}

\begin{Definition}{Степенной ряд}{}
	Степенным рядом с центром в точке $z_0 \in \mathbb C$ называется ряд
	\[
	\sum_{n=0}^{\infty} a_n (z-z_0)^n,\qquad a_n \in \mathbb C.
	\]
	Число $R \ge 0$ называется \emph{радиусом сходимости}, если ряд сходится при $|z-z_0|<R$ и расходится при $|z-z_0|>R$.
\end{Definition}

\begin{Theorem}{Формула Коши--Адамара}{}
	Для степенного ряда $\sum_{n=0}^{\infty} a_n (z-z_0)^n$ радиус сходимости равен
	\[
	R = \frac{1}{\limsup\limits_{n\to\infty}|a_n|^{1/n}},
	\]
	где $1/0=+\infty$, $1/(+\infty)=0$.
\end{Theorem}

\begin{proof}
	Обозначим $b_n:=|a_n|^{1/n}$. Без ограничения общности $z_0=0$.
	Ключевое наблюдение:
	\[
	|a_n z^n| = \bigl(|a_n|^{1/n}|z|\bigr)^n = (b_n|z|)^n.
	\]
	Рассмотрим три случая.
	
	\textbf{Случай 1: $\{b_n\}$ не ограничена.}
	Для любого $z\ne0$ найдётся подпоследовательность $\{n_k\}$ с $b_{n_k}>1/|z|$,
	откуда $(b_{n_k}|z|)^{n_k}>1$, и общий член ряда не стремится к нулю.
	Ряд расходится при всех $z\ne0$, то есть $R=0$.
	
	\textbf{Случай 2: $\{b_n\}$ ограничена, $L:=\limsup_{n\to\infty}b_n>0$.}
	Полагаем $R=1/L$.
	
	\emph{Сходимость при $|z|<R$.} Из $|z|<1/L$ следует $L<1/|z|$, поэтому найдётся
	$\varepsilon>0$ такое, что $L+\varepsilon<1/|z|$. По определению верхнего предела
	начиная с некоторого $N_0$ выполнено $b_n\le L+\varepsilon$, то есть
	\[
	b_n|z| \le (L+\varepsilon)|z| < 1-|z|\varepsilon =: q < 1.
	\]
	Следовательно, $|a_nz^n|=(b_n|z|)^n\le q^n$, и ряд сходится абсолютно.
	
	\emph{Расходимость при $|z|>R$.} Из $|z|>1/L$ следует $\limsup b_n\cdot|z|>1$,
	значит существует подпоследовательность $\{n_k\}$ с $(b_{n_k}|z|)^{n_k}>1$.
	Общий член ряда не стремится к нулю, ряд расходится.
	
	\textbf{Случай 3: $\lim_{n\to\infty}b_n=0$.}
	Для любого $z\in\mathbb C$ найдётся $N_0$ такое, что при $n\ge N_0$
	выполнено $b_n|z|<1/2$, откуда $|a_nz^n|\le(1/2)^n$.
	Ряд сходится абсолютно при всех $z$, то есть $R=+\infty$.
\end{proof}

\begin{Remark}{Поведение на границе круга сходимости}{}
	Формула Коши--Адамара ничего не говорит о сходимости при $|z-z_0|=R$.
	На границе возможно всё: сходимость во всех точках, расходимость во всех или смешанное поведение.
\end{Remark}

\begin{Example}{Граничное поведение}{}
	Ряд $\displaystyle\sum_{n=1}^{\infty}\frac{z^n}{n}$ имеет $R=1$:
	при $z=1$ расходится, при $z=-1$ сходится.
	Ряд $\displaystyle\sum_{n=1}^{\infty}\frac{z^n}{n^2}$ тоже имеет $R=1$,
	но сходится абсолютно на всей окружности $|z|=1$.
\end{Example}

\subsection{Равномерная сходимость внутри круга сходимости}

\begin{Theorem}{Равномерная сходимость и непрерывность суммы}{}
	Пусть ряд $\sum_{n=0}^{\infty}a_nz^n$ имеет радиус сходимости $R>0$
	и $0<r<R$. Тогда ряд сходится абсолютно и равномерно на $|z|\le r$,
	а его сумма $S(z)$ равномерно непрерывна на $|z|\le r$.
\end{Theorem}

\begin{proof}
	Выберем $r_1$ так, что $r<r_1<R$. Поскольку ряд $\sum|a_n|r_1^n$ сходится,
	его члены ограничены: $|a_n|r_1^n\le C$ для некоторой константы $C$ и всех $n$.
	При $|z|\le r$ и $m\ge1$:
	\[
	\left|\sum_{k=n+1}^{n+m}a_kz^k\right|
	\le\sum_{k=n+1}^{n+m}|a_k|r_1^k\cdot\left(\frac{r}{r_1}\right)^k
	\le C\sum_{k=n+1}^{n+m}\left(\frac{r}{r_1}\right)^k.
	\]
	Так как $r/r_1<1$, хвост геометрического ряда стремится к нулю при $n\to\infty$
	равномерно по $m$. По критерию Коши ряд сходится абсолютно и равномерно на $|z|\le r$.
	
	Докажем равномерную непрерывность $S$ на $|z|\le r$.
	Зафиксируем $\varepsilon>0$. Из равномерной сходимости найдётся $N_0$ такое, что
	\[
	\sup_{|z|\le r}|S_{N_0}(z)-S(z)|<\frac\varepsilon3,
	\]
	где $S_n(z)=\sum_{k=0}^na_kz^k$. Полином $S_{N_0}$ равномерно непрерывен на
	замкнутом круге $|z|\le r$, поэтому найдётся $\delta>0$ такое, что для всех
	$z_1,z_2$ с $|z_1|,|z_2|\le r$ и $|z_1-z_2|<\delta$:
	\[
	|S_{N_0}(z_1)-S_{N_0}(z_2)|<\frac\varepsilon3.
	\]
	Тогда по неравенству треугольника:
	\[
	|S(z_1)-S(z_2)|\le|S(z_1)-S_{N_0}(z_1)|+|S_{N_0}(z_1)-S_{N_0}(z_2)|+|S_{N_0}(z_2)-S(z_2)|<\varepsilon.\qedhere
	\]
\end{proof}

\begin{Corollary}{Непрерывность суммы}{}
	Сумма степенного ряда $S(z)=\sum_{n=0}^{\infty}a_nz^n$ непрерывна в открытом круге $|z|<R$.
\end{Corollary}

\begin{proof}
	Для любой точки $z_*$ с $|z_*|<R$ выберем $r$ так, что $|z_*|<r<R$.
	По теореме $S$ равномерно непрерывна на $|z|\le r$, а значит, непрерывна в точке $z_*$.
\end{proof}

\subsection{Почленное дифференцирование степенного ряда}

\begin{Theorem}{Почленное интегрирование и дифференцирование}{}
	Пусть ряд $\sum_{n=0}^{\infty}a_n(x-x_0)^n$ имеет радиус сходимости $R>0$
	и $0<r<R$. Тогда на отрезке $[x_0-r,\,x_0+r]$ этот ряд можно почленно интегрировать
	и дифференцировать. В частности, сумма $S(x)=\sum_{n=0}^{\infty}a_n(x-x_0)^n$
	дифференцируема при $|x-x_0|<R$, причём
	\[
	S'(x)=\sum_{n=1}^{\infty}n a_n(x-x_0)^{n-1}.
	\]
\end{Theorem}

\begin{proof}
	Без ограничения общности $x_0=0$. Ряд производных $\sum_{n\ge0}a_{n+1}(n+1)x^n$
	получен почленным дифференцированием; покажем, что его радиус сходимости равен $R$.
	
	По формуле Коши--Адамара радиус ряда производных равен
	\[
	R_1=\frac{1}{\limsup_{n\to\infty}|a_{n+1}(n+1)|^{1/n}}.
	\]
	Так как $(n+1)^{1/n}\to1$, получаем
	$\limsup|a_{n+1}(n+1)|^{1/n}=\limsup|a_{n+1}|^{1/n}$.
	Далее, $|a_{n+1}|^{1/n}=\bigl(|a_{n+1}|^{1/(n+1)}\bigr)^{(n+1)/n}$, и при $n\to\infty$
	показатель $(n+1)/n\to1$, поэтому
	\[
	\limsup_{n\to\infty}|a_{n+1}|^{1/n}=\limsup_{n\to\infty}|a_{n+1}|^{1/(n+1)}=\limsup_{n\to\infty}|a_n|^{1/n},
	\]
	откуда $R_1=R$.
	
	Следовательно, ряд производных равномерно сходится на $[-r,r]$, и по теореме о почленном
	дифференцировании равномерно сходящейся последовательности дифференцируемых функций
	$S'(x)=\sum_{n=1}^\infty na_nx^{n-1}$.
\end{proof}

\begin{Corollary}{Бесконечная дифференцируемость. Коэффициенты Тейлора}{}
	Сумма степенного ряда бесконечно дифференцируема внутри интервала сходимости, причём
	\[
	a_k=\frac{S^{(k)}(x_0)}{k!}.
	\]
\end{Corollary}

\begin{proof}
	Применяя предыдущую теорему $k$ раз:
	\[
	S^{(k)}(x)=\sum_{n\ge k}a_n\cdot n(n-1)\cdots(n-k+1)\cdot(x-x_0)^{n-k}.
	\]
	При $x=x_0$ все слагаемые с $n>k$ обнуляются, и $S^{(k)}(x_0)=a_k\cdot k!$.
\end{proof}

\subsection{Ряд Тейлора}

\begin{Definition}{Ряд Тейлора}{}
	Если $S$ бесконечно дифференцируема в точке $x_0$ и задаётся степенным рядом
	в её окрестности, то этот ряд есть \emph{ряд Тейлора} функции $S$:
	\[
	\sum_{k=0}^{\infty}\frac{S^{(k)}(x_0)}{k!}(x-x_0)^k.
	\]
\end{Definition}

\begin{Example}{Гладкая функция, не равная своему ряду Тейлора}{}
	Функция
	\[
	f(x)=\begin{cases}e^{-1/x^2},&x\ne0,\\0,&x=0\end{cases}
	\]
	бесконечно дифференцируема на $\mathbb R$, и все её производные в нуле равны нулю.
	Ряд Тейлора тождественно равен нулю, хотя $f(x)\ne0$ при $x\ne0$.
\end{Example}

\subsection{Бесконечные произведения}

\begin{Definition}{Бесконечное произведение}{}
	\emph{Бесконечным произведением} называется предел частичных произведений:
	\[
	\prod_{k=1}^{\infty}b_k:=\lim_{N\to\infty}\prod_{k=1}^{N}b_k.
	\]
	Произведение называется \emph{сходящимся}, если предел существует, конечен и отличен от нуля;
	иначе --- \emph{расходящимся}.
\end{Definition}

\begin{Statement}{Критерий через логарифмы}{}
	Пусть $b_k>0$ начиная с некоторого номера $n$. Тогда
	\[
	\prod_{k=1}^{\infty}b_k\;\text{ сходится}
	\iff
	\sum_{k=n}^{\infty}\ln b_k\;\text{ сходится}.
	\]
\end{Statement}

\begin{proof}
	Так как $\ln\prod_{k=1}^{N}b_k=\sum_{k=1}^{N}\ln b_k$ и логарифм непрерывен и строго монотонен
	на $(0,+\infty)$, существование конечного ненулевого предела произведения равносильно
	существованию конечного предела суммы логарифмов.
\end{proof}

\begin{Definition}{Абсолютная сходимость произведения}{}
	Произведение $\prod_{k=1}^\infty b_k$ называется \emph{абсолютно сходящимся},
	если $b_k>0$ начиная с некоторого номера и ряд $\sum_k|\ln b_k|$ сходится.
\end{Definition}

\begin{Corollary}{Произведения вида $\prod(1+a_k)$}{}
	Пусть $a_k>-1$ для всех $k$. Тогда произведение $\prod_{k=1}^{\infty}(1+a_k)$
	абсолютно сходится тогда и только тогда, когда сходится ряд $\sum_{k=1}^{\infty}|a_k|$.
\end{Corollary}

\begin{proof}
	($\Rightarrow$) Из абсолютной сходимости произведения следует сходимость $\sum|\ln(1+a_k)|$,
	откуда $a_k\to0$. При достаточно больших $k$ выполнено $|a_k|<1/2$, и для $|t|<1/2$:
	\[
	\frac{1}{2}\le\left|\frac{\ln(1+t)}{t}\right|\le\frac{3}{2}.
	\]
	Поэтому ряды $\sum|\ln(1+a_k)|$ и $\sum|a_k|$ (с некоторого номера) сходятся или расходятся
	одновременно.
	
	($\Leftarrow$) Из $\sum|a_k|<\infty$ следует $a_k\to0$, и та же оценка даёт сходимость
	$\sum|\ln(1+a_k)|$, то есть абсолютную сходимость произведения.
\end{proof}

\begin{Example}{Произведение для $\sin x$}{}
	\[
	\sin x = x\prod_{k=1}^{\infty}\!\left(1-\frac{x^2}{\pi^2k^2}\right),\qquad x\in\mathbb R.
	\]
\end{Example}

\begin{Example}{Гамма-функция}{}
	Формула Вейерштрасса:
	\[
	\Gamma(s)=\frac{1}{se^{\gamma s}}\prod_{k=1}^{\infty}\!\left(1+\frac{s}{k}\right)^{-1}e^{s/k},
	\qquad s\ne0,-1,-2,\ldots,
	\]
	где $\gamma=\lim_{n\to\infty}\!\bigl(\sum_{k=1}^{n}\tfrac1k-\ln n\bigr)$ --- постоянная Эйлера.
	Эквивалентное интегральное представление:
	\[
	\Gamma(s)=\int_0^{\infty}x^{s-1}e^{-x}\,dx.
	\]
\end{Example}

\subsection{Функции нескольких переменных}

Пространство $\mathbb R^d$ снабжается евклидовой нормой
$|x-y|=\sqrt{\sum_{k=1}^{d}(x_k-y_k)^2}$.
\emph{Областью} называется связное открытое множество $\Omega\subset\mathbb R^d$.

\begin{Definition}{Предел и непрерывность}{}
	Пусть $f:\Omega\subset\mathbb R^d\to\mathbb R$, $x$ --- предельная точка $\Omega$, $x\in\Omega$.
	Функция $f$ \emph{непрерывна} в точке $x$, если
	\[
	\lim_{\substack{y\to x\\y\in\Omega}}f(y)=f(x).
	\]
\end{Definition}

\begin{Definition}{Проверка по путям}{}
	Пусть $x$ --- внутренняя точка $\Omega$, $y\in\mathbb R^d$. \emph{Путём} в направлении $y$ называется функция
	\[
	\psi_y(t)=f(x+ty),\qquad t:\,x+ty\in\Omega.
	\]
	Если $f$ непрерывна в точке $x$, то для любого $y$ функция $\psi_y$ непрерывна в нуле.
	
	Для координатных направлений $y_k=(0,\ldots,1,\ldots,0)$
	это даёт непрерывность по каждой переменной в отдельности.
\end{Definition}

\begin{Remark}{Ограниченность метода путей}{}
	Непрерывность не равносильна непрерывности вдоль каждой координатной оси
	и даже не равносильна непрерывности вдоль каждого прямолинейного пути:
	могут понадобиться криволинейные пути.
\end{Remark}

\begin{Example}{Предел по прямым не гарантирует общего предела}{}
	Рассмотрим
	\[
	f(x_1,x_2)=\frac{x_1 x_2}{x_1^2+x_2^2},\qquad (x_1,x_2)\ne(0,0).
	\]
	По осям координат предел равен $0$. Вдоль прямой $x_1=x_2=t$ получаем $f(t,t)=1/2$.
	Следовательно, общего предела в нуле нет.
\end{Example}

\begin{Example}{Все прямолинейные пути дают ноль, но общего предела нет}{}
	Рассмотрим
	\[
	f(x_1,x_2)=
	\begin{cases}
		\dfrac{x_1 x_2^2}{x_1^2+x_2^4}, & (x_1,x_2)\ne(0,0),\\[6pt]
		0, & (x_1,x_2)=(0,0).
	\end{cases}
	\]
	Для любого фиксированного направления $y=(y_1,y_2)$:
	\[
	f(ty_1,ty_2)=\frac{ty_1\cdot t^2y_2^2}{t^2y_1^2+t^4y_2^4}\xrightarrow{t\to0}0
	\quad(y_1\ne0),
	\]
	то есть по всем прямолинейным путям предел равен $0$.
	Однако вдоль параболы $x_1=x_2^2$:
	\[
	f(t^2,t)=\frac{t^4}{2t^4}=\frac12,
	\]
	поэтому общего предела в нуле не существует.
\end{Example}

\subsection{Непрерывность отображений в $\mathbb{R}^n$}

\begin{Theorem}{Непрерывность на компакте}{}
	Если $f$ непрерывна на компакте $K\subset\mathbb{R}^d$, то $f$ равномерно непрерывна на $K$.
\end{Theorem}

Далее рассматриваем отображение
\[
f:\Omega\subset\mathbb{R}^d\to\mathbb{R}^n,\qquad
f(x)=\begin{pmatrix}f_1(x_1,\ldots,x_d)\\\vdots\\f_n(x_1,\ldots,x_d)\end{pmatrix}.
\]

Напомним: $x$ называется \emph{точкой сгущения} $\Omega$, если в любой окрестности $U(x)$ содержится бесконечно много точек $\Omega$, то есть $\#(U(x)\cap\Omega)=\infty$.

\begin{Statement}{Покомпонентная непрерывность}{}
	Пусть $f:\Omega\subset\mathbb{R}^d\to\mathbb{R}^n$ и $x$ --- точка сгущения $\Omega$. Тогда
	\[
	f\text{ непрерывна в }x
	\iff
	f_k\text{ непрерывна в }x\;\;\forall k=1,\ldots,n.
	\]
\end{Statement}

\begin{proof}
	($\Rightarrow$) Для каждой координаты:
	\[
	|f_k(x)-f_k(y)|\le\sqrt{\sum_{j=1}^n|f_j(x)-f_j(y)|^2}=|f(x)-f(y)|,
	\]
	поэтому из $|f(x)-f(y)|\to0$ следует $|f_k(x)-f_k(y)|\to0$.
	
	($\Leftarrow$) Для данного $\varepsilon>0$ из непрерывности каждой $f_k$ найдём $\delta_k>0$ такое, что при $0<|y-x|<\delta_k$ выполнено $|f_k(x)-f_k(y)|<\varepsilon$. Полагаем $\delta=\min_{1\le k\le n}\delta_k$. При $0<|y-x|<\delta$:
	\[
	|f(x)-f(y)|=\sqrt{\sum_{j=1}^n|f_j(x)-f_j(y)|^2}\le\sqrt{n\varepsilon^2}=\sqrt{n}\,\varepsilon.
	\]
	Это доказывает непрерывность $f$ в точке $x$.
\end{proof}

\begin{Theorem}{Непрерывность композиции}{}
	Пусть $f:\Omega\subset\mathbb{R}^d\to\mathbb{R}^n$, $g:\widetilde\Omega\subset\mathbb{R}^n\to\mathbb{R}^m$,
	$\widetilde\Omega\supset f(\Omega)$. Если $f$ непрерывна в точке $x$, а $g$ непрерывна в точке $f(x)$,
	то $g\circ f$ непрерывна в точке $x$.
\end{Theorem}

\begin{proof}
	Из $y\to x$ следует $f(y)\to f(x)$ (непрерывность $f$), а затем $g(f(y))\to g(f(x))$ (непрерывность $g$).
\end{proof}

\subsection{Дифференцируемость отображений}

\begin{Definition}{Дифференцируемость}{}
	Пусть $\Omega\subset\mathbb{R}^d$ --- открытое множество, $x\in\Omega$, $f:\Omega\to\mathbb{R}^n$.
	Функция $f$ называется \emph{дифференцируемой} в точке $x$, если существует линейный оператор
	$d_x:\mathbb{R}^d\to\mathbb{R}^n$ такой, что при $\xi\to x$, $\xi\in\Omega$:
	\[
	f(\xi)=f(x)+d_x(\xi-x)+o(|\xi-x|).
	\]
	Оператор $d_x$ единственен и называется \emph{дифференциалом} функции $f$ в точке $x$.
\end{Definition}

\begin{Remark}{Связь с одномерным случаем}{}
	При $d=n=1$ имеем $d_x=f'(x)$, и формула принимает стандартный вид
	$f(\xi)=f(x)+f'(x)(\xi-x)+o(|\xi-x|)$.
\end{Remark}

\begin{Statement}{Дифференцируемость влечёт непрерывность}{}
	Если $f:\Omega\to\mathbb{R}^n$ дифференцируема в точке $x$, то $f$ непрерывна в точке $x$.
\end{Statement}

\begin{proof}
	Из $f(\xi)-f(x)=d_x(\xi-x)+o(|\xi-x|)$ и ограниченности $d_x$:
	\[
	|f(\xi)-f(x)|\le\|d_x\|\cdot|\xi-x|+o(|\xi-x|)\to0\quad\text{при }\xi\to x.
	\]
\end{proof}

\begin{Statement}{Покомпонентная дифференцируемость}{}
	Если $f:\Omega\to\mathbb{R}^n$ дифференцируема в точке $x$, то каждая координатная функция
	$f_k:\Omega\to\mathbb{R}$ дифференцируема в точке $x$, причём дифференциалом $f_k$ служит
	линейный функционал
	\[
	d_{x,k}(h):=\bigl(d_x(h)\bigr)_k.
	\]
\end{Statement}

\begin{proof}
	Переходя к $k$-й координате в равенстве
	$f(\xi)=f(x)+d_x(\xi-x)+o(|\xi-x|)$:
	\[
	f_k(\xi)=f_k(x)+\bigl(d_x(\xi-x)\bigr)_k+o(|\xi-x|).
	\]
	Отображение $h\mapsto(d_x h)_k$ линейно, значит $f_k$ дифференцируема с $d_{x,k}=(d_x)_k$.
\end{proof}

\begin{Statement}{Частные производные дифференцируемой функции}{}
	Пусть $f:\Omega\subset\mathbb{R}^d\to\mathbb{R}$ дифференцируема в точке $x$. Тогда для каждого
	$j=1,\ldots,d$ существует частная производная
	\[
	\frac{\partial f}{\partial x_j}(x)
	:=\lim_{\theta\to0}\frac{f(x+e_j\theta)-f(x)}{\theta}
	=d_x(e_j),
	\]
	где $e_j=(0,\ldots,1,\ldots,0)$ --- $j$-й базисный вектор. Более того,
	\[
	d_x(\xi-x)=\sum_{j=1}^d\frac{\partial f}{\partial x_j}(x)\cdot(\xi_j-x_j).
	\]
\end{Statement}

\begin{proof}
	Зафиксируем $j$ и возьмём $\xi:=x+e_j\theta$ при достаточно малом $\theta$ (так что $\xi\in\Omega$).
	Тогда $|\xi-x|=|\theta|$, и из дифференцируемости:
	\[
	f(x+e_j\theta)=f(x)+d_x(e_j)\cdot\theta+o(|\theta|).
	\]
	Разделив на $\theta$ и перейдя к пределу:
	\[
	\frac{\partial f}{\partial x_j}(x)
	=\lim_{\theta\to0}\frac{f(x+e_j\theta)-f(x)}{\theta}
	=d_x(e_j).
	\]
	Таким образом, $j$-я компонента функционала $d_x$ равна $\partial f/\partial x_j(x)$.
	Раскладывая $\xi-x=\sum_{j=1}^d e_j(\xi_j-x_j)$ и используя линейность $d_x$:
	\begin{align*}
		f(\xi)
		&=f(x)+d_x(\xi-x)+o(|\xi-x|)\\
		&=f(x)+\sum_{j=1}^d d_x(e_j)\cdot(\xi_j-x_j)+o(|\xi-x|)\\
		&=f(x)+\sum_{j=1}^d\frac{\partial f}{\partial x_j}(x)\cdot(\xi_j-x_j)+o(|\xi-x|).\qedhere
	\end{align*}
\end{proof}

\subsection{Дифференцирование композиции}

\begin{Theorem}{Производная сложной функции}{}
	Пусть $f:\Omega\subset\mathbb{R}^d\to\mathbb{R}^n$, $g:\widetilde\Omega\subset\mathbb{R}^n\to\mathbb{R}^m$,
	$f(\Omega)\subset\widetilde\Omega$. Если $f$ дифференцируема в точке $x$, а $g$ дифференцируема
	в точке $f(x)$, то $g\circ f$ дифференцируема в точке $x$, причём
	\[
	d_x(g\circ f)=d_{f(x)}g\circ d_xf.
	\]
\end{Theorem}

\begin{proof}
	Пусть $\xi\to x$. Из дифференцируемости $f$:
	\[
	f(\xi)=f(x)+d_xf(\xi-x)+o(|\xi-x|),
	\]
	откуда $f(\xi)\to f(x)$ и $|f(\xi)-f(x)|=O(|\xi-x|)$.
	Из дифференцируемости $g$ в точке $f(x)$:
	\[
	g(f(\xi))=g(f(x))+d_{f(x)}g\bigl(f(\xi)-f(x)\bigr)+o\bigl(|f(\xi)-f(x)|\bigr).
	\]
	Подставляя $f(\xi)-f(x)=d_xf(\xi-x)+o(|\xi-x|)$ и используя $|f(\xi)-f(x)|=O(|\xi-x|)$:
	\[
	g(f(\xi))=g(f(x))+(d_{f(x)}g\circ d_xf)(\xi-x)+o(|\xi-x|).\qedhere
	\]
\end{proof}